\section{Tomboy / bokukko / prince-like girl: гендерная игра, сила, прямота, защитница}

Tomboy/bokukko/prince-like girl — амплуа гендерной игры: персонаж, который занимает социальную позицию, обычно закреплённую за мужской ролью — защитника, того, кто действует первым, — не как маскировку, а как честную идентичность. Утена Тенджоу из "Revolutionary Girl Utena" — образцовый случай именно потому, что сериал прямо исследует, что стоит за желанием "стать принцем", а не просто носить эту роль как костюм: детская травма, обещание, данное в момент отчаяния, и попытка жить внутри этого обещания годы спустя.

Механика барьера здесь не про тепло, как у dere-типов из главы 3, а про allowed action: этому амплуа разрешено то, что обычно требует оправдания у более "женственных" ролей — вступиться первой, отказаться защищать в ответ, если не хочется. Гендерная игра тут не декорация, а снятие ограничения на инициативу.

Баг этого амплуа предсказуем и стоит того, чтобы быть названным: тот, кто всегда защищает, систематически не умеет быть защищённым — принимать помощь, признавать, что сам нуждается в опоре, читается им как поражение роли, а не просто как неудобство. У Утены это не мелочь для колорита, а центральный механизм трагедии сериала.

Для Рин это амплуа даёт полезный, но не буквальный урок: её инициатива (2.7, "персонаж как двигатель") не должна требовать оправдания или разрешения — она вступается, потому что видит нужный момент, а не потому что кто-то её попросил. Но если довести эту инициативность до предела без противовеса, получится персонаж, который структурно не умеет принять заботу в свой адрес, — а Рин это уметь обязана, иначе её собственная субъектность (0.1, 1.4) останется только декларацией.
